\homeworktitle{Домашняя работа}{Логарифм: в какую степень?}
\studentline

% В polyglossia нет \asbuk (это babel-овский счётчик), поэтому определяем
% его здесь — кириллические буквы для enumerate-подпунктов.
\makeatletter
\def\asbuk#1{\@asbuk{\csname c@#1\endcsname}}
\def\@asbuk#1{%
  \ifcase#1\or
    а\or б\or в\or г\or д\or е\or ж\or з\or и\or к\or л\or м\or
    н\or о\or п\or р\or с\or т\or у\or ф\or х\or ц\or ч\or ш\or щ\or
    э\or ю\or я\else ??\fi
}
\makeatother

\begin{routebox}[Ключевая связь]
  \[
    \log_a b=c
    \quad\Longleftrightarrow\quad
    a^c=b,
    \qquad
    a>0,\quad a\ne1,\quad b>0.
  \]
  Логарифм отвечает на вопрос: «В какую степень нужно возвести число \(a\),
  чтобы получить число \(b\)?»
\end{routebox}

\begin{routebox}[Маршрут домашней работы]
  \begin{enumerate}[itemsep=0.2em,topsep=0.25em]
    \item Восстановите степени и смысл новой записи.
    \item Переведите равенства из одной формы в другую.
    \item Вычислите логарифмы через вопрос «в какую степень?».
    \item Решите простейшие уравнения.
    \item Объясните ошибку и выполните перенос метода.
  \end{enumerate}
  Задания 1--9 составляют обязательную часть. Задание 10* выполняется по
  желанию.
\end{routebox}

\section*{Часть I. Смысл логарифма}

\begin{homeworktask}[Найдите показатели степеней]{1}
  Заполните пропуски. Пункт г) проверяет, умеете ли вы работать с дробным
  показателем степени.
  \[
    \text{а) }2^{\mathgap[6mm]}=64;
    \qquad
    \text{б) }10^{\mathgap[6mm]}=0{,}001;
  \]
  \[
    \text{в) }4^{\mathgap[6mm]}=\frac1{16};
    \qquad
    \text{г) }81^{\mathgap[6mm]}=3.
  \]
\end{homeworktask}

\begin{homeworktask}[Переведите в логарифмическую форму]{2}
  \begin{enumerate}[itemsep=0.45em]
    \item[\textbf{а)}] \(3^4=81\);
    \item[\textbf{б)}] \(7^0=1\);
    \item[\textbf{в)}] \(5^{-2}=\dfrac1{25}\);
    \item[\textbf{г)}] \(16^{\frac34}=8\).
  \end{enumerate}
\end{homeworktask}
Для пункта г) заполните карту перевода. Заполняйте её по порядку: определите
основание, показатель и значение степени, а затем запишите равносильное
логарифмическое равенство. Остальные пункты оформите без карты.
\logtranslationmap

\begin{homeworktask}[Переведите в степенную форму]{3}
  \begin{enumerate}[itemsep=0.45em]
    \item[\textbf{а)}] \(\log_2 128=7\);
    \item[\textbf{б)}] \(\log_9 1=0\);
    \item[\textbf{в)}] \(\log_3\dfrac1{81}=-4\);
    \item[\textbf{г)}] \(\log_{25}125=\dfrac32\).
  \end{enumerate}
\end{homeworktask}
\answerspace{28mm}

\newpage
\section*{Часть II. Вычисление логарифмов}

\begin{homeworktask}[Вычислите логарифмы]{4}
  \[
    \text{а) }\log_2 64;
    \quad
    \text{б) }\log_7 1;
    \quad
    \text{в) }\log_5\frac1{125};
  \]
  \[
    \text{г) }\log_{16}2;
    \quad
    \text{д) }\log_{\frac13}27;
    \quad
    \text{е) }\lg0{,}0001.
  \]

  Для пункта \textbf{г)} заполните карту рассуждения. Пункт \textbf{д)}
  решите аналогичным способом без карты.
\end{homeworktask}
\logcalculationmap

\section*{Часть III. Решение уравнений}

\begin{homeworktask}[Запишите точный ответ]{5}
  Решите уравнения. Если значение показателя не удаётся определить с помощью
  известных степеней, запишите точный ответ в логарифмической форме.
  \[
    \text{а) }2^x=32;
    \qquad
    \text{б) }3^x=7;
    \qquad
    \text{в) }10^x=0{,}2.
  \]
  \[
    \text{г) }4^x=\frac18.
  \]
\end{homeworktask}
\answerspace{30mm}

\begin{homeworktask}[Перейдите к степенной записи]{6}
  Решите уравнения:
  \[
    \text{а) }\log_3 x=4;
    \qquad
    \text{б) }\log_2(x-5)=3;
    \qquad
    \text{в) }\log_5(2x+1)=2.
  \]
  \[
    \text{г) }\log_2(x-1)=-2.
  \]
  После решения проверьте, что аргумент каждого логарифма положителен.
\end{homeworktask}
\answerspace{52mm}

\begin{homeworktask}[Когда логарифм существует?]{7}
  \begin{enumerate}[itemsep=0.45em]
    \item[\textbf{а)}] Укажите, какие выражения имеют смысл:
      \[
        \log_2 8,
        \qquad
        \log_{-2}8,
        \qquad
        \log_1 5,
        \qquad
        \log_3(-9),
        \qquad
        \log_{\frac12}4.
      \]
    \item[\textbf{б)}] Найдите основание \(a\), если
      \[
        \log_a16=4.
      \]
      Учтите ограничения на основание логарифма.
  \end{enumerate}
\end{homeworktask}
\answerspace{36mm}

\newpage
\section*{Часть IV. Объясняем и переносим метод}

\begin{homeworktask}[Найдите и исправьте ошибку]{8}
  Ученик написал:
  \[
    \log_8 64=8,
    \qquad\text{потому что}\qquad 8\cdot8=64.
  \]
  Объясните, в чём ошибка. Запишите вопрос, на который отвечает данный
  логарифм, и исправьте решение.
\end{homeworktask}
\answerspace{38mm}

\begin{homeworktask}[Сравните две записи]{9}
  Из равенства \(4^3=64\) составлены две записи:
  \[
    \log_4 64=3,
    \qquad
    \log_{64}4=\frac13.
  \]
  Докажите обе записи переходом к степенной форме. Объясните, почему при
  перестановке основания и аргумента изменилось значение логарифма.
\end{homeworktask}
\answerspace{43mm}

\Needspace{14\baselineskip}
\begin{homeworktask}[Задание со звёздочкой]{10*}
  Вычислите:
  \[
    \text{а) }\log_{\sqrt2}8;
    \qquad
    \text{б) }\log_{\frac1{27}}9;
    \qquad
    \text{в) }\log_{0{,}01}10.
  \]

  \begin{hintbox}
    Представьте основание и аргумент как степени одного и того же числа.
    Например, \(\sqrt2=2^{\frac12}\).
  \end{hintbox}
\end{homeworktask}
\answerspace{45mm}

\begin{selfcheck}
  Перед сдачей проверьте:
  \begin{itemize}[itemsep=0.18em,topsep=0.25em]
    \item в переводе записи основание степени осталось основанием логарифма;
    \item каждый вычисленный логарифм можно проверить равенством \(a^c=b\);
    \item отрицательные и дробные показатели не были отвергнуты;
    \item в уравнениях выполнен переход к степенной форме;
    \item ответы вида \(\log_a b\) оставлены точными, если подбор невозможен.
  \end{itemize}
\end{selfcheck}
